\documentclass[a4paper,12pt]{ujarticle}
\usepackage[top=25mm,bottom=25mm,left=25mm,right=25mm]{geometry}
\title{週間報告}
\author{金武俊佑}
\date{2025年4月28日}

\usepackage{otf}
\usepackage{url}
\usepackage{tabularx}
\usepackage{comment}
\usepackage[dvipdfmx]{graphicx}
\usepackage[dvipdfmx]{color}
\usepackage{listings}
\usepackage{pdfpages}
\usepackage{multirow}
\usepackage{float}

\newcommand{\thc}[1]{\multicolumn{1}{c|}{#1}}

\begin{document}
\maketitle

\section{今週の報告}\label{sec:thisweek}
今週の報告を以下に記す。
\subsection{AtCoder}\label{sec:atcoder}
土曜日はバイトがあったので月曜日の面接後に挑戦したA問題は正解できたがB問題は正解できなかった。
アカウント名はLastiara
\subsection{就活}\label{sec:job}
就活の進捗を以下に記す。


現在は6社選考に進んでいる。


抱えている選考が少し多い気がするので、企業の厳選をして研究に支障がない範囲に調整していきたい。
\subsection{卒論:インフォディオ共同研究(AIモデル担当)}\label{sec:thesis}
卒論の進捗を以下に記す。


火曜日にコミュニティ版のDifyの環境構築ができた。


月曜日の夕方にDifyの教科書を見ながらアプリケーション開発の勉強をした。Chapter3の内容をすべて終わらせたかったが、時間が足りず途中までとなった。もう少し進めたいのでこれから眠くなるまで勉強したい。
\subsection{その他}\label{sec:other}
最近体力がなく面接やゼミ終わりに一時間ほど仮眠をとるような状況が続いているので、散歩などして体力をつけていきたい。
\section{来週の予定}\label{sec:nextweek}
来週の活動予定を以下に記す。
\begin{itemize}
  \item 04/30(木) 13:00 miniPBLに参加する。
\end{itemize}


\section{スケジュール}\label{sec:schedule}
Q1のスケジュールを表\ref{tab:studyschedule}に示す。
年間スケジュールを表\ref{tab:yearlyschedule}に示す。

\begin{table}[H]
  \caption{Q１のスケジュール}\label{tab:studyschedule}
  \centering
  \begin{tabular}{|p{4mm}||p{30mm}|p{30mm}|p{30mm}|p{30mm}|p{30mm}|}\hline
        & \thc{月} & \thc{火} & \thc{水} & \thc{木} & \thc{金} \\ \hline\hline
     1  &     &     &     &  &   \\ \hline
     2  &     &  &  &  & もくもく会\\ \hline
     3  & もくもく会  & ゼミ &  & &もくもく会  \\ \hline
     4  &  & ゼミ & &  & \\ \hline
  \end{tabular}
\end{table}

\begin{table}[H]
  \caption{年間スケジュール}\label{tab:yearlyschedule}
  \centering
  \begin{tabular}{|p{30mm}||p{30mm}|}\hline
        & \thc{内容} \\ \hline\hline
     4～5月  & テーマ詳細化・先行研究調査 \\ \hline
     6～7月  & システム設計 \\ \hline
     8～10月  & 実装 \\ \hline
     11月  & 実験・評価 \\ \hline
     11～12月  & 論文執筆（並行して修正） \\ \hline
     2027年1月  & 最終修正・提出・発表 \\ \hline
  \end{tabular}
\end{table}


\end{document}
